% !TEX program = xelatex
\documentclass[11pt,a4paper]{ctexart}

\usepackage{amsmath,amssymb,amsfonts}
\usepackage{booktabs}
\usepackage{geometry}
\usepackage{hyperref}
\usepackage{enumitem}
\usepackage{xcolor}
\usepackage{array}
\usepackage{listings}

\geometry{left=2.5cm,right=2.5cm,top=2.5cm,bottom=2.5cm}
\hypersetup{colorlinks=true,linkcolor=blue,urlcolor=blue}

\lstset{
  basicstyle=\ttfamily\small,
  breaklines=true,
  frame=single,
  columns=fullflexible,
  keepspaces=true
}

\newcolumntype{L}[1]{>{\raggedright\arraybackslash}p{#1}}

\title{%
  \textbf{星闪 SLB 物理层}\\[0.4em]
  \Large 附录 L 双向复合信道计算方法技术文档\\[0.3em]
  \large IQ 相乘复合 \textbullet\ 相位消偏 \textbullet\ 距离解算封装
}
\author{依据 T/XS 10001-2025《星闪无线通信系统 接入层\\
同步低时延宽带空口 SLB 技术要求和测试方法》\\
附录 L（资料性）整理并工程化\\
（团体标准 20250326 版）\\[0.3em]
\small 配套 6.6.2.3 基于双向复合信道的测距；给出公式解读、处理步骤与函数封装}
\date{2026 年 7 月 21 日}

\begin{document}
\maketitle
\tableofcontents
\newpage

%======================================================================
\part{总览}
%======================================================================

\section{文档范围与模块划分}

本文档对应 T/XS 10001-2025 \textbf{附录 L}（资料性）``被测设备数据包参数双向复合信道的计算方法''，说明 I/R 两端 CSI 如何合成为往返等效信道，以及如何据此估计距离。工程上服务于 \textbf{6.6.2.3} 解算节点的``双向复合 $\rightarrow$ 拼接 $\rightarrow$ 解距''链路。

\begin{table}[h]
\centering
\small
\caption{文档模块划分}
\begin{tabular}{clp{8cm}}
\toprule
\textbf{块} & \textbf{主题} & \textbf{核心输出} \\
\midrule
A & 协议对应与边界 & 附录 L 与 6.6.2.3 / 6.5.5.6 的关系 \\
B1 & 标准公式提炼 & 两端 CSI、相位、复合相位与 $r$ 的关系 \\
B2 & 处理步骤与封装 & 函数分层、伪代码、接口表 \\
C & 约束与衔接 & CFO/SCO、实现注意、章节指针 \\
\bottomrule
\end{tabular}
\end{table}

\section{与相关章节的关系}

\begin{table}[h]
\centering
\small
\begin{tabular}{lll}
\toprule
\textbf{节号} & \textbf{关系} & \textbf{说明} \\
\midrule
附录 L & 本文主体 & 资料性；给出 IQ 相乘消偏方法 \\
6.6.2.3.1 & 上游流程 & 规定要对 $N\cdot M$ 双向信道做复合并解距 \\
6.6.2.3.3 & 上游数据 & CSI 反馈汇聚后再调用本方法 \\
6.5.5.6 & 上游量定义 & 反馈 CSI（IQ）的测量量语义 \\
6.6.4 & 下游 & 距离 $D$ 可选融合为坐标 \\
\bottomrule
\end{tabular}
\end{table}

\section{方法在解算链中的位置}

\begin{enumerate}[nosep]
  \item 空口：两端按 6.6.2.3 采集并反馈各载波 CSI；
  \item \textbf{本方法}：对对齐后的两端 IQ 做双向复合，得到仅与距离相关的相位；
  \item 跨子载波 / 跨跳频载波合并处理，估计距离 $r$（或 $D$）。
\end{enumerate}

\newpage
%======================================================================
\part{附录 L 双向复合信道计算}
%======================================================================

\section{协议章节对应}

对应 T/XS 10001-2025 \textbf{附录 L}（资料性）。正文 6.6.2.3 要求``根据复合信道状态信息解算距离''，具体 IQ 相乘与消偏关系以本附录为计算参考。

\section{功能定义}

\begin{enumerate}[nosep]
  \item 根据 I/R 两端在子载波 $k$ 上的 CSI，建立含传播距离、时钟差、本振相差的相位模型；
  \item 通过对两端 IQ 相乘得到双向复合信道，消去时钟偏移与 LO 相差；
  \item 使复合相位仅与目标距离 $r$ 相关，从而可跨频率合并解距；
  \item 要求对 $f_k^{R}-f_k^{I}$ 做 CFO/SCO 等精确补偿。
\end{enumerate}

\section{模块边界（上下游及职责划分）}

\subsection{上游模块}

\begin{table}[h]
\centering
\small
\begin{tabular}{L{3.5cm}L{4.5cm}L{5cm}}
\toprule
\textbf{上游} & \textbf{输入} & \textbf{说明} \\
\midrule
6.6.2.3 测量与反馈 & 两端对齐的 CSI IQ & 按子载波 / 端口对 \\
6.5.5.6 & CSI 量化与语义 & IQ 复数值 \\
CFO/SCO 估计与补偿 & 补偿后的等效频率 & 附录 L 强制要求 \\
\bottomrule
\end{tabular}
\end{table}

\subsection{下游模块}

\begin{table}[h]
\centering
\small
\begin{tabular}{L{3.5cm}L{4.5cm}L{5cm}}
\toprule
\textbf{下游} & \textbf{输出} & \textbf{说明} \\
\midrule
距离估计器 & 距离 $r$ 或 $D$ & 相位斜率 / CIR 峰等 \\
6.6.4 & 坐标 & 可选 \\
\bottomrule
\end{tabular}
\end{table}

\subsection{本模块不负责的内容}

\begin{itemize}[nosep]
  \item 跳频状态机、空口发收与 CSI 反馈组装（见 6.6.2.3）；
  \item CSI 信道估计算法本身（接收机实现）；
  \item 坐标解算（见 6.6.4）。
\end{itemize}

\section{在 SLB 通信流程中的角色}

位于\textbf{解算节点}侧、CSI 汇聚完成之后：消费第一/第二位置节点上报的 CSI，产出复合信道与距离，供定位应用或 6.6.4 使用。

%======================================================================
\section{关键参数与强制要求}
%======================================================================

\begin{table}[h]
\centering
\small
\begin{tabular}{L{3cm}L{5.5cm}L{4.5cm}}
\toprule
\textbf{量} & \textbf{含义} & \textbf{约束 / 来源} \\
\midrule
$h_k^{R}(t_1)$ / $h_k^{I}(t_2)$ & R/I 节点在子载波 $k$ 的 CSI & 附录 L \\
$\Delta t$ & I--R 时钟偏差 & 复合后应被消去 \\
$\theta_{IR}$ & 双方 LO 相位差 & 复合后应被消去 \\
$r$ & 目标距离 & 复合相位仅保留与 $r$ 相关项 \\
$c_0$ & 光速 & 距离换算 \\
$f_k^{R}-f_k^{I}$ & 两端在子载波 $k$ 的频率差 & \textbf{宜精确补偿}（CFO/SCO） \\
\bottomrule
\end{tabular}
\end{table}

%======================================================================
\section{附录 L 核心详解}
%======================================================================

\subsection{两端 CSI 模型}

I 节点与 R 节点做信道估计，第 $k$ 个子载波 CSI 为：
\begin{align}
  h_k^{R}(t_1)
  &= h_{k,0}\,
    e^{-j\bigl[2\pi(f_k^{R}-f_k^{I})t_1
    +\varphi_{I,\mathrm{init}}-\varphi_{R,\mathrm{init}}\bigr]}, \\
  h_k^{I}(t_2)
  &= h_{k,0}\,
    e^{-j\bigl[2\pi(f_k^{I}-f_k^{R})t_2
    +\varphi_{R,\mathrm{init}}-\varphi_{I,\mathrm{init}}\bigr]}.
\end{align}

\subsection{单端相位（含距离、时钟、LO）}

记 $\Delta t$ 为 clock offset，$\theta_{IR}$ 为 LO phase offset，则：
\begin{align}
  \varphi_R(f_k,r)
  &= -2\pi(f_k^{R}-f_k^{I})
     \Bigl(\frac{r}{c_0}+\Delta t\Bigr)
     -\theta_{IR}
     \pmod{2\pi}, \\
  \varphi_I(f_k,r)
  &= -2\pi(f_k^{I}-f_k^{R})
     \Bigl(\frac{r}{c_0}-\Delta t\Bigr)
     +\theta_{IR}
     \pmod{2\pi}.
\end{align}

\subsection{IQ 相乘：双向复合}

对对应 IQ 值相乘 $h_k^{R}(t)\,h_k^{I}(t)$，相位相加后消去 $\Delta t$ 与 $\theta_{IR}$：
\begin{equation}
  \label{eq:phi2w}
  \varphi_{2W}(f_k,r)
  =\varphi_R(f_k,r)+\varphi_I(f_k,r)
  =-4\pi(f_k^{R}-f_k^{I})\,\frac{r}{c_0}.
\end{equation}

\textbf{工程结论}：
\begin{itemize}[nosep]
  \item 复合后相位\textbf{仅与目标距离 $r$ 相关}；
  \item 不同 $f_k$ 上的 $\varphi_{2W}$ \textbf{可合并处理}以解 $r$；
  \item $f_k^{R}-f_k^{I}$ 宜精确补偿，包括载波频偏（CFO）与采样时钟偏差（SCO）。
\end{itemize}

在已补偿、两端工作于同一绝对频率网格 $f_k$ 的实现中，常将往返相位写为等效形式
$\varphi_{2W}(f_k)\approx-4\pi f_k\,r/c_0$，从而
\begin{equation}
  r=-\frac{c_0}{4\pi}\frac{\mathrm{d}\varphi_{2W}}{\mathrm{d}f}.
\end{equation}

%======================================================================
\section{数据类型、长度与维度}
%======================================================================

\begin{table}[h]
\centering
\small
\begin{tabular}{llll}
\toprule
\textbf{对象} & \textbf{类型/长度} & \textbf{单位} & \textbf{备注} \\
\midrule
单 tone CSI & 复数 IQ & --- & 每子载波 1 个 \\
复合 CSI $H^{2W}(k)$ & 复数 & --- & $H_R\cdot H_I$ \\
相位 $\varphi_{2W}$ & float & rad & 需 unwrap \\
距离 $r$ & float & m & 由斜率或 CIR 峰得到 \\
频率轴 $f_k$ & float & Hz & 补偿后绝对/差频轴 \\
\bottomrule
\end{tabular}
\end{table}

%======================================================================
\section{时序关系与触发条件}
%======================================================================

\begin{itemize}[nosep]
  \item \textbf{触发}：解算节点收齐某一载波（或已拼接频段）上双方对齐的 CSI；
  \item \textbf{先决}：子载波索引一致；建议已完成 CFO/SCO 补偿；
  \item \textbf{之后}：跨 $k$ 合并解距；跳频场景下先 per-carrier 复合再跨载波拼接。
\end{itemize}

%======================================================================
\section{处理步骤}
%======================================================================

\subsection{推荐分层（airMode = slb）}

\begin{table}[h]
\centering
\small
\begin{tabular}{L{5cm}L{8cm}}
\toprule
\textbf{函数} & \textbf{职责} \\
\midrule
\texttt{slbCompensateCfoScoForCsi} & （可选前置）补偿 $f_k^{R}-f_k^{I}$ 相关频偏 \\
\texttt{slbCalculateCompositeCsi} & 两端 IQ 逐 tone 相乘 $\rightarrow H^{2W}$ \\
\texttt{slbEstimateDistanceFromCompositeCsi} & 由 $\varphi_{2W}(f)$ 或 IFFT 估 $r$ \\
\texttt{slbResolveCompositeRange} & 对外入口：对齐 $\rightarrow$ 复合 $\rightarrow$ 解距 \\
\bottomrule
\end{tabular}
\end{table}

\subsection{Step 流水线}

\begin{enumerate}[nosep]
  \item \textbf{Step 1 对齐}：按子载波编号（全反馈 160 或表 18 部分反馈）对齐 $H^{R}[k]$、$H^{I}[k]$；
  \item \textbf{Step 2 补偿}：对 CFO/SCO 做精确补偿（附录 L 要求）；
  \item \textbf{Step 3 复合}：对每个 $k$，$H^{2W}[k]=H^{R}[k]\cdot H^{I}[k]$（复数乘法）；
  \item \textbf{Step 4 取相}：$\varphi[k]=\mathrm{atan2}(\Im H^{2W},\Re H^{2W})$，并相位解缠绕；
  \item \textbf{Step 5 解距}：对 $\varphi(f)$ 拟合斜率，或 IFFT 找往返峰，换算 $r$；
  \item \textbf{Step 6（跳频）}：多载波重复 Step 1--3 后拼接再 Step 4--5，或每载波解距后融合。
\end{enumerate}

\subsection{复合内核（实现）}

\begin{lstlisting}[language=C,caption={slbCalculateCompositeCsi 内核}]
/* H2W = HR * HI  (附录 L：IQ 相乘) */
for (k = 0; k < toneCount; ++k) {
    out[k].re = hr[k].re * hi[k].re - hr[k].im * hi[k].im;
    out[k].im = hr[k].re * hi[k].im + hr[k].im * hi[k].re;
}
\end{lstlisting}

\subsection{距离估计（与式 \eqref{eq:phi2w} 衔接）}

\textbf{方法 A — 相位斜率（直接对应附录 L``可合并处理''）}
\begin{enumerate}[nosep]
  \item 计算并 unwrap $\varphi_{2W}(k)$；
  \item 对频率轴拟合 $\varphi\approx s\cdot f+\varphi_0$；
  \item $r=-(c_0/(4\pi))\,s$（在等效 $\varphi\approx-4\pi f r/c_0$ 约定下）。
\end{enumerate}

\textbf{方法 B — IFFT / CIR}
\begin{enumerate}[nosep]
  \item 对 $H^{2W}(f)$ 做 IFFT 得 $|h(\tau)|$；
  \item 取主峰往返时延 $\tau_{\mathrm{rt}}$；
  \item $r=c_0\cdot\tau_{\mathrm{rt}}/2$。
\end{enumerate}

\subsection{模块接口表}

\begin{table}[h]
\centering
\small
\begin{tabular}{llll}
\toprule
\textbf{方向} & \textbf{信号/参数} & \textbf{类型} & \textbf{描述} \\
\midrule
in & \texttt{csiNodeR[]} & complex & R 端各 tone CSI \\
in & \texttt{csiNodeI[]} & complex & I 端各 tone CSI \\
in & \texttt{freqHz[]} & float & 补偿后频率轴 (Hz) \\
in & \texttt{toneCount} & uint & tone 数 \\
in & \texttt{speedOfLightMps} & float & $c_0$ (m/s) \\
out & \texttt{compositeCsi[]} & complex & $H^{2W}$ \\
out & \texttt{distanceM} & float & 估计距离 $r$ (m) \\
\bottomrule
\end{tabular}
\end{table}

%======================================================================
\section{约束与实现要点}
%======================================================================

\begin{enumerate}[nosep]
  \item \textbf{【资料性】}：附录 L 为资料性；6.6.2.3 正文只要求复合后解距，实现应以本附录公式为默认参考。
  \item \textbf{【IQ 相乘】}：复合操作为复数乘法；相位相加自动消去 $\Delta t$ 与 $\theta_{IR}$。
  \item \textbf{【CFO/SCO】}：$f_k^{R}-f_k^{I}$ 宜精确补偿，否则 $\varphi_{2W}$ 斜率失真，距离偏差。
  \item \textbf{【子载波对齐】}：两端 tone 集合必须一致（全反馈或同一 $P$ 的表 18 集合）。
  \item \textbf{【相位解缠】}：斜率法必须 unwrap，否则距离跳变。
  \item \textbf{【与 $T_3=T_4$】}：空口测量单元对称等待（6.6.2.3.1）保证两端 CSI 时刻关系符合本附录模型前提。
  \item \textbf{【跳频】}：先 per-carrier 复合，再跨载波拼接增大等效带宽，提高距离分辨率。
  \item \textbf{【端口维】}：对每个端口对 $(i,j)$ 独立调用复合+解距，再融合。
\end{enumerate}

%======================================================================
\section{与标准其他章节的衔接}
%======================================================================

\begin{itemize}[nosep]
  \item \textbf{附录 L}：本文公式与消偏方法来源；
  \item \textbf{6.6.2.3}：测量、反馈与``复合后解距''流程入口；
  \item \textbf{6.5.5.6}：CSI IQ 测量量定义；
  \item \textbf{6.6.4}：距离到坐标（可选）。
\end{itemize}

\vfill
\begin{center}
\small\textcolor{gray}{完整表述以 T/XS 10001-2025 附录 L 为准；本文档提供双向复合信道计算的工程解读、处理步骤与函数封装，供 6.6.2.3 解算实现与联调参考。}
\end{center}

\end{document}
